\documentclass[11pt,a4paper]{article}

% Packages
\usepackage[utf8]{inputenc}
\usepackage[T1]{fontenc}
\usepackage{amsmath,amssymb}
\usepackage{graphicx}
\usepackage{hyperref}
\usepackage{xcolor}
\usepackage{geometry}
\usepackage{setspace}
\usepackage{natbib}
\usepackage{abstract}

% Page geometry
\geometry{margin=1in}

% Hyperref settings
\hypersetup{
    colorlinks=true,
    linkcolor=blue,
    citecolor=blue,
    urlcolor=blue
}

% Title and author
\title{\textbf{Compression, Release, and Evolution: A Unified Systems Theory from Entropy}}
\author{
    田陆 (Lu Tian)\\
    \textit{Independent Researcher}\\
    \texttt{uglyboy@aliyun.com}
}
\date{First published: 2026-05-25\\White paper version: 2026-06-02}

\begin{document}

\maketitle

\begin{abstract}
\noindent Systems across physics, biology, and society maintain local order by redistributing entropy to other dimensions. We propose a framework anchored in a redefinition of control: not as a target state, but as an ongoing process that transports randomness (entropy) across dimensions via compression mapping ($f: X \rightarrow Y$, $|Y| < |X|$). From this redefinition, entropy becomes the natural descriptive language rather than an external metaphor. Antifragility emerges as an independent property, irreducible to rigidity or flexibility. Coding transforms system evolution into a unified analytical framework. The theory identifies three structural fallacies shared across eight disciplines---coding reversal, dimension compression bias, and self-referential closure---and argues that evolution is the precondition for stability, not merely a desirable feature. Applying the framework's logic to coding evolution itself reveals that the three-beat coding-evolution cycle is the universal paradigm of evolution, a structural consequence of the theory's own axioms. Formalization and empirical validation remain as next steps.
\end{abstract}

\vspace{0.5cm}

\section{Introduction}

The second law of thermodynamics states that total entropy in an isolated system never decreases. Yet organisms grow, technologies iterate, and civilizations accumulate complexity. This apparent contradiction between the thermodynamic arrow and the arrow of complexification has been a central puzzle since Boltzmann.

Existing resolutions fall into two traditions. Prigogine's dissipative structures and Schrödinger's ``negative entropy'' explain order as the product of open systems that export entropy to their environment \citep{prigogine1984, schrodinger1944}. Holling's adaptive cycle and panarchy framework describe how complex systems cycle through growth, conservation, release, and reorganization \citep{holling1973, gunderson2002}. Both capture essential aspects of the puzzle, but neither provides a unified mechanism connecting the thermodynamic roots of order to observable failure modes across disciplines.

This paper proposes a theory whose core move is a redefinition of control itself. Rather than defining control by its endpoint---a target state reached or a constraint satisfied---we define it as an ongoing process: \textbf{control is the transportation of randomness across dimensions}. The operational mechanism is a \textbf{compression mapping} ($f: X \rightarrow Y$, where $|Y| < |X|$), and the natural language for describing what is transported is \textbf{entropy}, defined axiomatically as randomness. From this process-based definition, we derive antifragility as an independent system property. We argue that coding---the formal expression of control rules---transforms system evolution from domain-specific descriptions into a unified analytical framework. We conclude that evolution is not merely a desirable property of healthy systems; it is the prerequisite for stability itself.

From this mechanism, we develop a diagnostic framework for analyzing system failure. We identify three cross-disciplinary fallacy patterns and trace them across eight disciplines. We propose antifragility---the balance between rigidity (control intensity) and flexibility (strategy space)---as a cross-disciplinary meta-standard for evaluating theories, policies, and institutional designs.

Our approach complements traditional science. Conventional methods ask, ``Given initial conditions, what will happen?'' (state derivation). Our theory asks, ``Given the current distribution of control, which dimensions are being hollowed out by invisible entropy leakage under random shocks?'' (risk assessment). The theory finds strongest application in social, organizational, and institutional domains where control is shallow and leakage is high. It aligns with the second law of thermodynamics, draws inspiration from biological evolution and ecological resilience, and aims to derive novel diagnostic predictions for code-governed systems. The derivation across all claimed domains, however, remains a research program rather than an established result. The two paradigms converge where control is deep---particle physics, orbital mechanics---and diverge where it is shallow---social systems, economic policy.

\textbf{Roadmap.} Section 2 develops the core theoretical vocabulary: entropy, compression mapping, rigidity, flexibility, antifragility, and coding evolution. It builds toward the coding-evolution cycle and the universal paradigm of evolution (\S2.6). Section 3 situates the framework relative to existing theories. Section 4 applies the theory diagnostically, identifying three cross-disciplinary fallacy patterns across eight disciplines. Section 5 translates the theory into a practical diagnostic protocol. Section 6 discusses limitations, a path to formalization, and complementarity with traditional science.

\section{Theoretical Framework}

\subsection{Entropy}

A precise description of control requires a language for randomness. We adopt \textbf{entropy}, defined as \textbf{randomness}: the number of distinct states a system can occupy. This definition aligns with the thermodynamic form ($S = k \ln W$, where $W$ counts microstates) and the information-theoretic form ($H = -\sum p_i \ln p_i$). In all three, entropy quantifies the size of a possibility space.

A room with books scattered across the floor, glasses on random surfaces, and cables in arbitrary configurations has high entropy. The number of possible arrangements is combinatorially large. After tidying, with every object in a designated position, the room has low entropy in the object-position dimension: only one or a few arrangements correspond to ``tidy.''

\textbf{Technical and analogical uses of entropy.} This framework uses entropy in two registers. In physical and information-theoretic systems (\S3.1), entropy carries its technical definition (consistent with Boltzmann and Shannon) and can be quantified in well-defined state spaces. In social, organizational, and institutional systems, entropy serves as a heuristic concept whose formalization requires first establishing the relevant state-space definition (see \S6.3). Both registers share the same conceptual structure---measuring the size of a possibility space---but differ in their degree of operationalization. Throughout this paper, we signal the distinction explicitly when the application domain shifts.

\subsection{Control as Compression Mapping}

\textbf{The redefinition of control.} Existing theories define control through its outcome: a system reaches a target state or satisfies a constraint. In this view, control is measured by the \textit{state} that results. Machines keep temperature within a band; organizations hit quarterly targets; laws maintain social order. Wiener's cybernetics improved on static definitions by recognizing control as a dynamic process. However, the process it described was the \textit{evolution of the controller}---how feedback loops adapt over time---not the mechanism of control itself. What computations occur during control? What is being transformed, and into what? In the standard account, control remains a black box labeled ``restriction'' or ``regulation,'' described by its endpoint rather than its operation.

We propose a different definition: \textbf{control is the process of transporting randomness from one dimension to another}. It is not the final state of constraint; it is the ongoing act of relocation. The implementing mechanism is a \textbf{compression mapping}:

\begin{center}
\textbf{$f: X \rightarrow Y$, where $|Y| < |X|$}
\end{center}

Multiple states in the input space $X$ map to the same state in the output space $Y$. When a team is told to ``use Python for the backend,'' the programming-language dimension is compressed from dozens of candidates to one. When a recipe specifies ``use only three seasonings,'' the seasoning dimension is compressed from a large combinatorial space to a small one. The control is not the resulting homogeneity; it is the continuous operation of forcing many possibilities to converge on few.

\textbf{Why entropy emerges naturally.} Because control is a process that transports randomness, entropy becomes the natural descriptive language. If control merely produced a final state, entropy would be unnecessary. We could describe control entirely in terms of initial and target states. But once we ask \textit{what was done to the randomness during this process}, we need to track it. This redefinition does not borrow entropy as an external metaphor; it makes entropy inevitable for describing what control \textit{is}.

\textbf{The critical consequence:} total entropy does not decrease. The randomness removed from the compressed dimension does not vanish. It leaks into dimensions not covered by the mapping. Compressing the seasoning dimension forces more intense innovation in cooking technique, heat control, and ingredient pairing to compensate. Total cooking complexity has not diminished; it has redistributed.

This principle is isomorphic to the second law of thermodynamics: no process reduces total entropy. What appears as entropy reduction in a focal dimension is always accompanied by entropy increase elsewhere. \textbf{Control is entropy transportation, not entropy elimination.}

\subsection{Rigidity and Flexibility}

Compression mapping as the mechanism of control yields two poles on a continuous spectrum of control intensity:

\textbf{Rigidity} equals more control. A rigid system applies compression mapping to more dimensions, or applies tighter constraints within a given dimension. The total control quantity is larger, and the remaining strategy space---the set of alternative paths---is narrower. A glass bowl is an extreme case: almost every molecule is locked into a crystal lattice, leaving virtually no strategy space. When dropped, the bowl has exactly one path: fracture along crystal planes.

\textbf{Flexibility} equals less control. A flexible system applies fewer compression mappings, or applies looser constraints. The total control quantity is smaller, and the remaining strategy space is wider. A sponge can deform, bounce, and absorb impact because its structure leaves room for molecular rearrangement.

Rigidity and flexibility are \textbf{neutral descriptive terms}. Neither is inherently good or bad. A system is ``too rigid'' when its strategy space is too narrow to respond to events outside its compressed pathways. A system is ``too flexible'' when its control is too weak to maintain coherent structure; it ceases to be a system and becomes random drift.

\subsection{Antifragility}

We define \textbf{antifragility} as a qualitative indicator of how much randomness a system can handle in the dimensions it controls. It is not a precise scalar: the randomness-handling capacities of two different dimensions cannot be directly compared. Rather, antifragility provides a direction for reasoning: before any decision, the question is not ``Does this make a specific metric more efficient?'' but ``Does this increase the system's overall antifragility?''

\textbf{Antifragility is an independent property.} It is not reducible to rigidity, flexibility, or any simple combination of the two. Rigidity describes control intensity; flexibility describes strategy space width; antifragility describes the system's capacity to absorb and metabolize randomness while maintaining core function. These are three distinct axes. A rigid system may or may not be antifragile depending on how its control is distributed; a flexible system may or may not be antifragile depending on whether its freedom is structured. The practical consequence is that changing rigidity alone---either increasing or decreasing control---cannot guarantee improved antifragility. The diagnostic question is always: which dimensions are being controlled, and which are being starved of entropy-release capacity?

The relationship between antifragility, rigidity, and flexibility is directional rather than definitional:

\begin{itemize}
    \item \textbf{Too rigid}: strategy space is too narrow; the system collapses when a random event falls outside its narrow path. Examples include cancer cells (all control compressed into proliferation, killing the host) and glass bowls (zero strategy space for impact).
    \item \textbf{Too flexible}: control is too weak; the system cannot maintain coherent structure. Example: a team with no task allocation, no quality standards, and no coordination---each member has maximum freedom, but the system as a whole has collapsed.
\end{itemize}

\textbf{Antifragility is defined on the system as a whole, not on individual controls.} Multiple controls can operate simultaneously, but if they counteract each other without forming a coherent force, antifragility does not increase. The art of system design is not maximizing control; it is drawing the correct \textbf{dimension distribution map}: which dimensions need control, and which need room to breathe.

Three common misconceptions:

\begin{enumerate}
    \item Rigidity is not antifragility. More control narrows strategy space. A stone is extremely rigid but shatters on impact.
    \item Flexibility is not antifragility. Removing all control increases freedom but destroys coordination. Antifragility requires both control force and strategy space.
    \item Speed and pressure do not equal optimality. Acceleration increases entropy production; simultaneous control increase narrows the dissipation channels. The result is internal entropy accumulation and eventual burnout, analogous to placing an air conditioner's outdoor unit inside the room.
\end{enumerate}

\textbf{Operational form.} Within this framework, an antifragility assessment takes the comparative form: \textit{under shock X, does the system's core function survive, degrade, or improve?} The answer is directional (``more antifragile than'' rather than ``antifragility score of N'') and specific to the dimensions being shocked. An antifragility claim is disconfirmed if a system exposed to a specified class of random shocks in a specified dimension exhibits degradation of core function rather than maintenance or improvement.

\subsection{System, Subsystem, and Heat Dissipation}

A \textbf{system} is a set of elements operating under shared control, forming a coherent whole. It is not a static object but a continuous process of control maintenance. A band is not ``four people with instruments''; it is those four people continuously playing according to a shared score. When the playing stops, the system ceases to exist.

Every system must release entropy to its environment. Maintaining low entropy in certain dimensions---order---necessarily produces high entropy in others. If the entropy-production rate exceeds the \textbf{heat dissipation bandwidth}---the capacity of available dimensions to absorb and release entropy---the system accumulates entropy internally until it collapses. A system that runs faster produces more entropy and requires more dissipation bandwidth. Speed without bandwidth expansion makes a system more fragile, not stronger.

A \textbf{subsystem} exists within a parent system. In some dimensions, the subsystem maintains low entropy not through its own control but because the parent system provides it. A cell does not regulate its own temperature; the organism does. This arrangement is also the subsystem's vulnerability: if the parent system fails, the subsystem collapses in those dimensions.

\subsection{Coding and the Universal Paradigm of Evolution}

Sections 2.1--2.5 established what control is and how systems maintain it. They showed that control operates through compression mapping, the balance of rigidity and flexibility, and the management of heat dissipation. But systems do not merely maintain control; they also change. The question driving the remainder of the framework is: how do systems evolve, and is there a universal language for describing that evolution?

\textbf{Why coding.} To study how systems evolve, we need a common language for describing what evolves. A system's control rules---its regulations, syntax, genetic code---constitute its \textbf{coding}: the formal expression of control. An organization's regulations, a programming language's syntax, a species' genome: these are all codes. The philosophical idea that a formal symbolic system can capture the structure of a domain has a recognized precedent in Wittgenstein's picture theory of language \citep{wittgenstein1922}, which argued that propositions picture facts through shared logical form. In the present framework, however, codes are not mere representations; they are the operational rules that actively govern a system's entropy distribution.

The move from studying individual system evolutions to studying \textit{code evolutions} is the theory's key methodological contribution. Without coding, each system type requires its own evolutionary language. Biological evolution proceeds through mutation and selection, legal evolution through precedent and legislation, technological evolution through design and iteration. These operate with separate vocabularies, separate communities, and separate methods. Coding collapses this diversity into a single analytical framework: a genome mutation, a contract clause, and a software patch are all the same operation---a code modification that alters entropy distribution---and each can be analyzed with the same tools. This unification enables the cross-disciplinary analysis in Section 4 and makes the theory a usable diagnostic instrument rather than merely a philosophy of systems.

\textbf{The coding-evolution cycle.} Once systems are described through codes, their evolution follows a single positive-feedback dynamic:

\begin{quote}
\textbf{Entropy decrease (new codes) $\rightarrow$ Entropy increase (new-old combinations) $\rightarrow$ New entropy decrease $\rightarrow$ ...}
\end{quote}

The first half-beat is entropy decrease. A new code forms: a constraint that reduces randomness in a specific dimension of expression. Poetic form is itself a code---it governs which combinations of syllables, tones, and imagery are permissible---and the history of poetry is a history of codes evolving. The Tang Dynasty's regulated verse (\textit{jueju}) imposed strict syllabic and tonal constraints on poetry---a compression mapping on linguistic expression.

The second half-beat is entropy increase. The new code combines with existing codes, generating an explosion of combinatorial possibilities. The constraints of regulated verse, far from killing poetry, forced poets to explore imagery, inversion, and wordcraft in dimensions everyday language never needed. The information density of regulated verse far exceeds that of ordinary speech.

The cycle then repeats: new combinatorial richness exposes new dimensions that need control, prompting the formation of yet newer codes. Song Dynasty \textit{ci} poetry loosened some constraints; modern free verse loosened others. Each relaxation triggers a new round of entropy increase on the combinatorial dimension, followed by new constraints.

This cycle operates without a conscious ``coder.'' Gravity spontaneously compresses matter into galactic disks---a compression mapping without a designer. Human-designed codes (laws, programming languages) and spontaneously emerged codes (physical laws, market conventions) both instantiate the same mechanism.

Why does this cycle never stop? Every round of entropy decrease---gaining control---necessarily produces more entropy in the code-usage dimension. More entropy means more randomness, which creates new control demands---a positive feedback loop. In plain terms: every new rule creates situations the old rules did not cover, which forces the creation of yet more rules. There is no terminal state of ``perfect control''; the combinatorial space of code interactions grows factorially with the number of codes, far exceeding any control system's capacity to cover.

\textbf{Why the cycle is universal.} The three-beat form of the coding-evolution cycle is not merely an empirical observation; it is a structural consequence of the theory's own logic. When coding evolution is itself analyzed as a system, the ``dimension'' being operated on is the code-space itself---the set of all possible control rules. A compression mapping on this space takes the form of a new code: a constraint that reduces the entropy of expression in some sub-dimension.

Forming new codes is not merely \textit{one} way to control the code-space; it is the \textit{only} viable one. This follows from the premise that the code-space contains only codes and their interactions. Any control operation on this space---including a meta-level constraint on which codes are permitted---is itself describable as a code and therefore constitutes another compression mapping on the code-space. There is no operation available at this level that is not a code. The entropy displaced by this compression is released into the code-combination dimension, generating new code-code interactions. The three-beat cycle is therefore the universal paradigm of evolution, a structural consequence of the theory's own axioms.

\textbf{Evolution as a system.} The process of evolution can itself be analyzed through the theory's lens, with its own control mechanisms, strategy space, and antifragility. Natural selection is the rigid face: a compression mapping that eliminates unfit individuals, converging the population toward adaptive states. Genetic drift (random variation) is the flexible face: small-scale random perturbations that explore adjacent possibilities. Recombination operates at the code level, splicing code fragments to produce new combinations, which selection then filters. A healthy evolutionary system requires balance: excessive drift makes convergence too slow (too flexible); excessive selection pressure flattens diversity (too rigid); rich recombination with moderate selection yields both strategy space and filtering efficiency.

Evolution also evolves its own mechanisms. The transitions from asexual to sexual reproduction, from purely genetic transmission to cultural transmission (language), and from waiting for real mutations to counterfactual reasoning each represent a meta-level compression mapping on the ``means of evolution'' dimension. These create new layers of control that dramatically accelerate the evolutionary cycle.

If the coding-evolution paradigm is universal, then any complex adaptive system---including life itself---should be analyzable as an instance of it.

\textbf{Life as an instance of the paradigm.} Schrödinger famously described life as feeding on ``negative entropy,'' absorbing order from the environment to maintain internal order \citep{schrodinger1944}. Within the present framework, this can be reframed as life seeking entropy in specific dimensions---novel information, external stimuli, environmental energy---and metabolizing it to maintain low entropy in critical dimensions. This reframing offers a candidate explanation for novelty-seeking behavior, consistent with dopamine responses to novel stimuli \citep{schultz1998}, but remains a hypothesis. From the theory's perspective, life is not a special case; it is one instance of a universal coding-evolution cycle operating on a particular substrate.

\section{Relation to Existing Theories}

The proposed framework integrates and extends several established intellectual traditions. It does not seek to replace them; it provides a common language for identifying what they share.

\subsection{Thermodynamics and Information Theory}

The theory's definition of entropy as randomness is consistent with both thermodynamic entropy \citep{boltzmann1877, gibbs1878} and information-theoretic entropy \citep{shannon1948}. The shared mathematical structure---a functional that is additive, convex, and positive-definite, measuring the size of a possibility space---supports the extension to social, organizational, and institutional systems. In these contexts, the ``possibility space'' refers to the set of behaviors, strategies, or states available to agents.

\subsection{Cybernetics and Control Theory}

Ashby's Law of Requisite Variety \citep{ashby1956} states that a controller must possess at least as much variety as the system it controls: ``only variety can destroy variety.'' The proposed theory provides a mechanism-level complement to this principle. Ashby answered \textit{how much} control is needed; this theory answers \textit{how} control works, through compression mapping that transfers entropy across dimensions. The entropy-leakage concept explains why increasing control in one dimension often fails to improve overall system stability: entropy does not disappear; it migrates to dimensions outside the controller's monitoring scope.

Stafford Beer's Viable System Model \citep{beer1972, beer1979} and management cybernetics share the theory's concern with how organizations maintain viability through recursive control structures. The coding-evolution framework offers a complementary perspective: Beer focused on the structural architecture of viable systems; the present theory focuses on the dynamic process by which those architectures evolve.

\subsection{Ecological Resilience and Panarchy}

Holling's adaptive cycle \citep{holling1973} and the panarchy framework \citep{gunderson2002} describe how social-ecological systems cycle through four phases: exploitation (r), conservation (K), release ($\Omega$), and reorganization ($\alpha$). The conservation (K) phase corresponds to increasing rigidity: control accumulates, connections become rigid, and the system becomes vulnerable to shocks outside its narrowed strategy space. The release ($\Omega$) phase corresponds to entropy leakage reaching a critical threshold, triggering dimension collapse. The reorganization ($\alpha$) phase corresponds to the formation of new codes.

The key difference is explanatory ambition: the proposed theory provides a candidate mechanism---compression mapping followed by entropy leakage---that may explain \textit{why} the adaptive cycle takes the form it does, rather than merely describing \textit{what} occurs. This hypothesis requires formal demonstration. The theory also extends the domain beyond ecosystems to include legal systems, programming languages, organizational management, and other code-governed systems.

\subsection{Institutional Economics and Collective Action}

Ostrom's design principles for long-enduring common-pool resource institutions \citep{ostrom1990, ostrom2005} identify structural features---clear boundaries, low-cost conflict resolution, minimal recognition rights, nested governance---that correlate with institutional survival. In the proposed framework: clear boundaries define the code's domain; low-cost conflict resolution provides heat-dissipation channels; nested governance establishes multi-layer coding architecture. Ostrom's empirical finding that systems satisfying 7--8 principles have $>$85\% long-term survival rates, while those satisfying fewer than 3 all collapse \citep[ch. 3]{ostrom1990}, provides an empirical anchor for the claim that code structure determines system antifragility.

North's institutional change theory \citep{north1990} describes how institutions (codes) reduce uncertainty (specific-dimension entropy decrease) while creating new interest conflicts and strategic spaces (entropy leakage to other dimensions). The proposed theory's coding-evolution cycle formalizes this dynamic.

\subsection{Complexity Economics}

Arthur's complexity economics \citep{arthur2013, arthur2021} views markets as non-equilibrium, emergence-driven systems where agents continuously adapt their strategies. Beinhocker \citep{beinhocker2006} describes economic evolution as a search algorithm operating on a space of ``business plans'' (codes). The proposed theory aligns with these frameworks but adds a diagnostic dimension: it asks not only how emergence happens, but what kind of emergence is not fragile. It focuses on the dimension distribution of control rather than the fact of emergence itself.

\subsection{Novel Contributions}

While individual components of the theory have precedents, the following elements represent novel synthesis:

\begin{enumerate}
    \item The unified mechanism of \textbf{control as compression mapping} provides a common language for control operations across physical, biological, social, and institutional domains.
    
    \item The \textbf{cross-domain generalization of the entropy-leakage principle}---that the dynamic of entropy transfer across dimensions, well-documented in thermodynamics and ecology, applies with equal structural force to social and institutional systems---offers a testable diagnostic for why optimizing one dimension often degrades overall system health.
    
    \item The \textbf{three cross-disciplinary fallacy patterns} (coding reversal, dimension compression bias, self-referential closure) identify structural errors that recur across disciplines, offering a meta-level diagnostic that no single-discipline framework captures.
    
    \item The hypothesized inversion of \textbf{``life as entropy resister'' to ``life as entropy seeker''} provides a candidate foundation for understanding motivation, innovation, and the dynamics of cultural evolution. This idea is consistent with the framework but requires further empirical support (see \S2.6).
\end{enumerate}

\section{Cross-Disciplinary Analysis: Three Fallacy Patterns}

The theory identifies three structural fallacy patterns that recur across eight disciplines. Each stems from the same root cause: compression mapping as a dimension-specific operation.

\subsection{Coding Reversal}

\textbf{Definition:} Each discipline mistakes an operational code for the system's fundamental purpose. The means of measurement becomes the goal of operation.

\textbf{Examples:}

\begin{itemize}
    \item \textbf{Economics:} ``Individual utility maximization'' is taken as the purpose of markets, when the deeper purpose is maintaining a transaction system that enables continuous code alignment among participants.
    
    \item \textbf{Political Science:} ``One person, one vote'' is taken as the purpose of democracy, when the deeper purpose is enabling continuous code evolution in the social system.
    
    \item \textbf{Medicine:} ``Eliminate the pathological target'' is taken as the purpose of health, when health is a multi-dimensional entropy balance that requires tracking where intervention-induced entropy leaks.
    
    \item \textbf{Education:} ``Standardized test scores'' are taken as the purpose of learning, when learning is the formation of an individual's own evolutionary engine.
\end{itemize}

\textbf{Mechanism:} Compression mapping is a reduction operation. When a complex multi-dimensional reality is mapped to a single metric ($f$: all activities $\rightarrow$ \{one number\}), the metric becomes the only dimension that ``counts.'' Over time, the system optimizes for the metric rather than the underlying reality it was designed to proxy.

\subsection{Dimension Compression Bias}

\textbf{Definition:} Codes inherently favor measurable dimensions. When ``only what is measured gets managed'' becomes the default strategy, immeasurable dimensions (buffers, redundancy, empathy, growth) are systematically ignored.

\textbf{Examples:}

\begin{itemize}
    \item \textbf{Lean production:} Zero-inventory, just-in-time delivery maximizes efficiency in the measurable dimension of inventory turnover, but eliminates the buffer dimension that absorbs supply-chain shocks. The COVID-19 pandemic revealed the fragility of this optimization \citep{choi2023}.
    
    \item \textbf{Scientific evaluation:} Impact factors, citation counts, and publication numbers compress all research activity into countable dimensions. Negative results, replication studies, and paradigm-challenging work---essential for science's self-correcting mechanism---are filtered out of the ``publishable'' category.
    
    \item \textbf{GDP-centered development:} GDP measures the aggregate value of goods and services but remains blind to households on the ``poverty knife-edge'' (one emergency from destitution), the erosion of social trust, and the depletion of natural capital.
\end{itemize}

\textbf{Mechanism:} This is a structural feature of code-based governance, not a failure of individual cognition. Codes are finite descriptions of control rules; a finite description can only encode a finite number of dimensions. The unencoded dimensions do not disappear. They continue to exist and accumulate entropy, but they lose institutional legitimacy because the system has no language to process them.

\subsection{Self-Referential Closure}

\textbf{Definition:} All internal critiques within a discipline are resolved within the existing code framework. Economics adds constraints to utility functions; management replaces bad KPIs with better ones; education responds to the shortcomings of standardization with more standards.

\textbf{Examples:}

\begin{itemize}
    \item \textbf{Economics:} Behavioral economics demonstrates that humans deviate from rational-agent assumptions, but the framework remains ``deviation from rationality'' rather than questioning whether rationality is the right benchmark. Risk management explains post-hoc why lean production was fragile, but cannot diagnose the fragility before a disaster.
    
    \item \textbf{Political Science:} Arrow's impossibility theorem proves that no voting system can satisfy all reasonable axioms \citep{arrow1951}; public choice theory reveals rent-seeking in voting; deliberative democracy advocates supplement voting with deliberation. All critiques operate within the assumption that voting is the framework; they describe what is lost, but do not question whether the compression mapping itself is the problem.
    
    \item \textbf{Management:} Goodhart's Law states that when a measure becomes a target, it ceases to be a good measure \citep{goodhart1975}. The response is to design better KPIs, balanced scorecards, and OKRs, all of which remain codes operating in the measurable-dimension space.
\end{itemize}

\textbf{Mechanism:} A discipline's code framework is also its legitimacy foundation. To question the code is to question the discipline's basis for existence. Internal critiques therefore converge on improving the code rather than replacing it. This is not intellectual dishonesty; it is a structural property of institutionalized codes.

\section{Diagnostic Framework}

\subsection{Three Questions}

For any system---a team, a body, a company, a policy---the theory provides a diagnostic protocol of three questions:

\begin{enumerate}
    \item \textbf{Where is control applied?} Identify the compression mappings currently operating in the system. Which dimensions are being compressed?
    
    \item \textbf{Which dimensions are approaching their limits?} Entropy never disappears; it transfers. When you compress the attendance dimension (requiring everyone to clock in by 10 AM), where does the displaced entropy flow? Is it leaking into spontaneous communication, innovation, or employee well-being? Assess the remaining heat-dissipation capacity of each dimension.
    
    \item \textbf{Where is the entropy flowing next?} If a saturated dimension continues to receive entropy, where will it transfer the overflow? Digestive system failure does not cause entropy to explode in place; it transfers to the immune system. Legal department compression does not eliminate compliance risk; it transfers undetected risk to business operations. Tracing entropy flow paths draws a collapse roadmap.
\end{enumerate}

\subsection{PDCA: A Practical Methodology}

The Plan-Do-Check-Act (PDCA) cycle, originally developed by Shewhart and Deming for total quality management \citep{deming1986}, maps directly onto the theory's operational logic:

\begin{itemize}
    \item \textbf{Plan:} Select a dimension and establish a target state. Ensure the target retains sufficient strategy space; do not choose a direction already near its rigidity limit.
    
    \item \textbf{Do:} Apply a compression mapping to converge the affected elements toward the target.
    
    \item \textbf{Check:} Verify that the compression mapping is effective and that entropy has not leaked into dangerous dimensions. If the compression introduces new instability, the mapping needs adjustment.
    
    \item \textbf{Act:} If the compression is successful and stable, solidify it as part of the system's dynamic control rules.
\end{itemize}

PDCA is a micro-scale implementation of the coding-evolution cycle. It enables rapid dimension adjustment without triggering phase transitions. The system evolves through controlled, incremental compression-and-check cycles rather than waiting for environmental selection to force change.

\subsection{Chinese Medicine as a System-Practice Precedent}

\begin{quote}
\textit{This section describes a diagnostic analogy. It does not endorse or evaluate specific medical claims of Chinese medicine.}
\end{quote}

Chinese medicine provides an illustrative case of pre-modern system-level diagnostic practice. Without modern analytical instruments, it infers internal entropy distribution from surface signals: tongue coating, pulse characteristics, facial complexion. In this analogy, acupuncture can be understood as a calibrated micro-perturbation on specific meridian dimensions that redirects entropy flow to unblock congested pathways. Herbal medicine adjusts the rhythm of entropy transformation and release rather than supplementing specific chemical compounds.

That three practitioners may give three different diagnoses for the same patient is not a defect under this interpretation; it is an inherent property of multi-dimensional system intervention. Different diagnoses may describe different entry points into the same underlying entropy-distribution imbalance. The principle of ``treating disease before it manifests'' (\textit{zhiweibing})---detecting entropy-circulation deviations through surface signals before symptoms appear---is, from the system-theoretic perspective, a direct application of the diagnostic framework to human physiology.

\section{Discussion}

\subsection{What This Theory Provides}

The theory does not supply a ``correct code'' to replace all others; that would itself commit coding reversal. Instead, it provides three tools:

\begin{enumerate}
    \item \textbf{Dimension diagnosis before code operation.} Before any measurement or optimization, the framework asks: What are this system's core survival dimensions? Which are covered by existing codes, and which are neglected? Under what conditions will entropy leakage from neglected dimensions breach critical thresholds?
    
    \item \textbf{Codes as evolution objects, not a priori frameworks.} A discipline's current core codes---utility functions, GDP, blood-pressure standards, admission scores---are historical products, not cosmic constants. They should iterate with system-state changes. In practice, codes acquire institutional inertia because changing them means changing the discipline's legitimacy foundation.
    
    \item \textbf{Antifragility as a cross-disciplinary meta-standard.} The test of a theory or practice is not what it measures, predicts, or controls, but whether it can maintain core functions under unexpected shocks and obtain structured improvement from those shocks.
\end{enumerate}

\subsection{Limitations}

The theory, in its current form, has several significant limitations:

\begin{enumerate}
    \item \textbf{Qualitative nature.} The framework provides conceptual direction but lacks formalized mathematical definitions, measurement protocols, and quantitative predictions. The claim that entropy leaks from dimension A to dimension B cannot currently be operationalized with specific units or thresholds.
    
    \item \textbf{No predictive track record.} The diagnostic claims---for example, that a system with compression mapping applied to its dissipation dimensions will collapse---have not been tested in a prospective, falsifiable study. The case analyses are retrospective.
    
    \item \textbf{Scope ambiguity.} The theory claims to cover physics, biology, and social systems under one framework, but the concept of ``entropy'' operates at different levels of abstraction in each domain. The bridging argument---that the mathematical structure is isomorphic across domains---has not been formally demonstrated.
    
    \item \textbf{No empirical calibration.} Unlike Ostrom's design principles (tested against 91+ case studies) or ecological resilience theory (which has identified measurable regime-shift thresholds), the theory has not been calibrated against empirical data.
\end{enumerate}

\subsection{Path to Formalization}

The most promising path to formalization builds on three empirical anchors in the literature:

\begin{itemize}
    \item \textbf{Ecological resilience theory} (Holling, Gunderson) provides measurable ``controlling variables''---the slow variables that determine system behavior, which are the ecological equivalent of the theory's ``dimensions.'' Documented regime shifts---for example, the shift from clear-water to turbid states in shallow lakes \citep{scheffer2001}, coral reef degradation, and rangeland desertification---provide measurable thresholds corresponding to the theory's ``entropy-leakage critical points.'' More recent work generalizing Ashby's law to multi-scale systems \citep{siegenfeld2025} suggests a formal pathway for extending requisite-variety reasoning across hierarchical levels.
    
    \item \textbf{Ostrom's design principles} \citep{ostrom1990} provide an operationalized checklist of code-structure properties that correlate with institutional survival. These may be reinterpreted as ``code antifragility indicators'' and tested against the theory's predictions.
    
    \item \textbf{Innis and McLuhan's media theory} \citep{innis1951, mcluhan1964} provides a generative mechanism for why codes become self-locking: the physical structure of the code medium biases what content is easily produced, transmitted, and institutionalized. This explains the mechanism behind dimension-compression bias.
\end{itemize}

A formalization program could proceed by (1) defining a multi-dimensional state space for a target system; (2) operationalizing compression mapping as variance reduction in selected dimensions; (3) defining antifragility as the probability of survival under multi-dimensional random shocks; and (4) testing the prediction that compression applied to dissipation dimensions reduces antifragility more than equivalent compression applied to control dimensions.

A minimal prospective test of the framework could take the following form. Identify a system currently undergoing compression in a measurable dimension---for instance, a regulatory tightening that constrains a specific industry practice. Using the three diagnostic questions (\S5.1), predict which adjacent dimension is most likely to absorb the displaced entropy---for instance, informal compliance workarounds, unreported risk externalization, or deterioration in an adjacent service quality. Specify observable indicators that would confirm the leakage prediction (e.g., a measurable increase in workaround behaviors within a defined time window) and indicators that would disconfirm it. A single such prospective study---predicting leakage before it is observed, rather than explaining it post-hoc---would move the framework from retrospective coherence to testable hypothesis.

\subsection{Complementarity with Traditional Science}

The theory's relationship to traditional science is one of complementarity, not replacement. In domains where control is deep and entropy leakage is minimal---particle physics, orbital mechanics, chemical kinetics---state-derivation methods are highly effective. The theory predicts and explains this: when control is near-complete and leakage is near-zero, risk assessment and state derivation should, in principle, yield equivalent conclusions. In domains where control is shallow and leakage is high---social systems, organizational design, economic policy---state-derivation methods systematically fail because they do not track entropy transfer across unmodeled dimensions. The theory provides a complementary diagnostic tool for these domains.

\section{Conclusion}

This paper began by redefining a foundational concept: control is not the achievement of a target state but the ongoing process of transporting randomness across dimensions. This redefinition makes entropy the natural language of control, reveals antifragility as an independent property rather than a blend of rigidity and flexibility, and---through the introduction of coding---transforms the study of system evolution from a collection of domain-specific narratives into a unified analytical framework.

The theory identifies three cross-disciplinary fallacy patterns---coding reversal, dimension compression bias, and self-referential closure---that recur across eight disciplines. These patterns are structural consequences of code-based governance, not the result of individual cognitive bias.

\textbf{The closed loop: evolution as the foundation of stability.} The most consequential result of treating evolution as a system (\S2.6) is the argument that a system's capacity to evolve is not a secondary consideration; it is the precondition for stability itself. A system that cannot evolve cannot sustain its structure under the entropy leakage produced by its own operations. The accumulation of entropy in unmanaged dimensions guarantees eventual collapse unless the system can modify its codes---that is, evolve. Stability is therefore not a property that can be designed once and preserved. It is a property that must be continuously regenerated through evolutionary motion. The diagnostic language this theory provides---compression mapping, entropy leakage, heat dissipation bandwidth, dimension distribution---is ultimately a language for assessing whether a system's evolutionary capacity is keeping pace with its entropy production.

The theory is currently qualitative. Its next steps involve operationalizing compression mapping, entropy leakage, and antifragility within measurable multi-dimensional state spaces, and testing the prediction that systems with insufficient evolutionary bandwidth deteriorate even when their momentary control metrics appear healthy. Until formalized, the theory offers a diagnostic coherence that no single-discipline framework provides: a common language for reasoning about which dimensions need control, which need release, and whether a system's evolutionary engine is running fast enough to outrun its own entropy.

\bibliographystyle{plainnat}

\begin{thebibliography}{26}

\bibitem[Arrow(1951)]{arrow1951}
Arrow, K.~J. (1951).
\newblock {\em Social Choice and Individual Values}.
\newblock Wiley.

\bibitem[Arthur(2013)]{arthur2013}
Arthur, W.~B. (2013).
\newblock {\em Complexity Economics: A Different Framework for Economic Thought}.
\newblock SFI Working Paper.

\bibitem[Arthur(2021)]{arthur2021}
Arthur, W.~B. (2021).
\newblock Foundations of complexity economics.
\newblock {\em Nature Reviews Physics}, 3, 136--145.

\bibitem[Ashby(1956)]{ashby1956}
Ashby, W.~R. (1956).
\newblock {\em An Introduction to Cybernetics}.
\newblock Chapman \& Hall.

\bibitem[Beer(1972)]{beer1972}
Beer, S. (1972).
\newblock {\em Brain of the Firm}.
\newblock Allen Lane.

\bibitem[Beer(1979)]{beer1979}
Beer, S. (1979).
\newblock {\em The Heart of Enterprise}.
\newblock Wiley.

\bibitem[Beinhocker(2006)]{beinhocker2006}
Beinhocker, E.~D. (2006).
\newblock {\em The Origin of Wealth}.
\newblock Harvard Business School Press.

\bibitem[Boltzmann(1877)]{boltzmann1877}
Boltzmann, L. (1877).
\newblock {\"Uber die Beziehung zwischen dem zweiten Hauptsatze der mechanischen W{\"armetheorie und der Wahrscheinlichkeitsrechnung}.
\newblock {\em Wiener Berichte}, 76, 373--435.

\bibitem[Choi et~al.(2023)]{choi2023}
Choi, T.~Y., Netland, T.~H., Sanders, N., Sodhi, M.~S., \& Wagner, S.~M. (2023).
\newblock Just-in-time for supply chains in turbulent times.
\newblock {\em Production and Operations Management}, 32(7), 2331--2340.

\bibitem[Deming(1986)]{deming1986}
Deming, W.~E. (1986).
\newblock {\em Out of the Crisis}.
\newblock MIT Center for Advanced Engineering Study.

\bibitem[Goodhart(1975)]{goodhart1975}
Goodhart, C. (1975).
\newblock Problems of monetary management: The U.K. experience.
\newblock {\em Papers in Monetary Economics}, 1, 1--20.

\bibitem[Gunderson and Holling(2002)]{gunderson2002}
Gunderson, L.~H. \& Holling, C.~S. (Eds.) (2002).
\newblock {\em Panarchy: Understanding Transformations in Human and Natural Systems}.
\newblock Island Press.

\bibitem[Holling(1973)]{holling1973}
Holling, C.~S. (1973).
\newblock Resilience and stability of ecological systems.
\newblock {\em Annual Review of Ecology and Systematics}, 4, 1--23.

\bibitem[Innis(1951)]{innis1951}
Innis, H.~A. (1951).
\newblock {\em The Bias of Communication}.
\newblock University of Toronto Press.

\bibitem[McLuhan(1964)]{mcluhan1964}
McLuhan, M. (1964).
\newblock {\em Understanding Media: The Extensions of Man}.
\newblock McGraw-Hill.

\bibitem[North(1990)]{north1990}
North, D.~C. (1990).
\newblock {\em Institutions, Institutional Change and Economic Performance}.
\newblock Cambridge University Press.

\bibitem[Ostrom(1990)]{ostrom1990}
Ostrom, E. (1990).
\newblock {\em Governing the Commons: The Evolution of Institutions for Collective Action}.
\newblock Cambridge University Press.

\bibitem[Ostrom(2005)]{ostrom2005}
Ostrom, E. (2005).
\newblock {\em Understanding Institutional Diversity}.
\newblock Princeton University Press.

\bibitem[Prigogine and Stengers(1984)]{prigogine1984}
Prigogine, I. \& Stengers, I. (1984).
\newblock {\em Order out of Chaos: Man's New Dialogue with Nature}.
\newblock Bantam Books.

\bibitem[Scheffer et~al.(2001)]{scheffer2001}
Scheffer, M., Carpenter, S., Foley, J.~A., Folke, C., \& Walker, B. (2001).
\newblock Catastrophic shifts in ecosystems.
\newblock {\em Nature}, 413(6856), 591--596.

\bibitem[Schr{\"odinger(1944)]{schrodinger1944}
Schr{\"odinger, E. (1944).
\newblock {\em What is Life?}
\newblock Cambridge University Press.

\bibitem[Schultz(1998)]{schultz1998}
Schultz, W. (1998).
\newblock Predictive reward signal of dopamine neurons.
\newblock {\em Journal of Neurophysiology}, 80(1), 1--27.

\bibitem[Shannon(1948)]{shannon1948}
Shannon, C.~E. (1948).
\newblock A mathematical theory of communication.
\newblock {\em Bell System Technical Journal}, 27, 379--423, 623--656.

\bibitem[Siegenfeld and Bar-Yam(2025)]{siegenfeld2025}
Siegenfeld, A.~F. \& Bar-Yam, Y. (2025).
\newblock A formal definition of scale-dependent complexity and the multi-scale law of requisite variety.
\newblock {\em Entropy}, 27(8), 835.

\bibitem[Taleb(2012)]{taleb2012}
Taleb, N.~N. (2012).
\newblock {\em Antifragile: Things That Gain from Disorder}.
\newblock Random House.

\bibitem[Wittgenstein(1922)]{wittgenstein1922}
Wittgenstein, L. (1922).
\newblock {\em Tractatus Logico-Philosophicus}.
\newblock (C.~K. Ogden, Trans.). Routledge \& Kegan Paul.

\end{thebibliography}

\vspace{1cm}

\begin{center}
\textit{This white paper is released under CC BY-NC-SA 4.0.}\\
\textit{First published 2026-05-25.}\\
\textit{For discussions, please open an issue at the project repository.}
\end{center}

\end{document}